\begin{lemma}
\label{lemma-essential-surjection}
Let $f: B \to A$ be a ring map in $\mathcal{C}_\Lambda$.
Notation as in (\ref{equation-sequence-functorial}).
\begin{enumerate}
\item The equivalent conditions of
Lemma \ref{lemma-surjective}
characterizing when $f$ is surjective are also equivalent to
\begin{enumerate}
\item $\Im(\text{d}_B) \to \Im(\text{d}_A)$ is surjective, and
\item the map $\Omega_{B/\Lambda} \otimes_B k \to
\Omega_{A/\Lambda} \otimes_A k$ is surjective.
\end{enumerate}
\item The following are equivalent
\begin{enumerate}
\item $f$ is an essential surjection
(see Lemma \ref{lemma-essential-surjection-mod-squares}),
\item the map $\Im(\text{d}_B) \to \Im(\text{d}_A)$ is an
isomorphism, and
\item the map $\Omega_{B/\Lambda} \otimes_B k \to
\Omega_{A/\Lambda} \otimes_A k$ is an isomorphism.
\end{enumerate}
\item If $k/k'$ is separable, then $f$ is an essential surjection if
and only if the map
$\mathfrak m_B/(\mathfrak m_\Lambda B + \mathfrak m_B^2) \to
\mathfrak m_A/(\mathfrak m_\Lambda A + \mathfrak m_A^2)$
is an isomorphism.
\item If $f$ is a small extension, then $f$ is not essential if and only if
$f$ has a section $s: A \to B$ in $\mathcal{C}_\Lambda$
with $f \circ s = \text{id}_A$.
\end{enumerate}
\end{lemma}

\begin{proof}
Proof of (1). It follows from (\ref{equation-sequence-functorial})
that (1)(a) and (1)(b) are equivalent. Also, if
$A \to B$ is surjective, then (1)(a) and (1)(b) hold. Assume (1)(a).
Since the kernel of $\text{d}_A$ is the image of
$H_1(L_{k/\Lambda})$ which also maps to
$\mathfrak m_B/\mathfrak m_B^2$ we conclude that
$\mathfrak m_B/\mathfrak m_B^2 \to \mathfrak m_A/\mathfrak m_A^2$
is surjective. Hence $B \to A$ is surjective by
Lemma \ref{lemma-surjective}. This finishes the proof of (1).

\medskip\noindent
Proof of (2). The equivalence of (2)(b) and (2)(c) is immediate from
(\ref{equation-sequence-functorial}).

\medskip\noindent
Assume (2)(b). Let $g : C \to B$ be a ring map in $\mathcal{C}_\Lambda$
such that $f \circ g$ is surjective. We conclude that
$\mathfrak m_C/\mathfrak m_C^2 \to \mathfrak m_A/\mathfrak m_A^2$
is surjective by
Lemma \ref{lemma-surjective}.
Hence
$\Im(\text{d}_C) \to \Im(\text{d}_A)$ is surjective
and by the assumption we see that
$\Im(\text{d}_C) \to \Im(\text{d}_B)$ is surjective.
It follows that $C \to B$ is surjective by (1).

\medskip\noindent
Assume (2)(a). Then $f$ is surjective and we see that
$\Omega_{B/\Lambda} \otimes_B k \to \Omega_{A/\Lambda} \otimes_A k$
is surjective. Let $K$ be the kernel. Note that
$K = \text{d}_B(\Ker(\mathfrak m_B/\mathfrak m_B^2 \to
\mathfrak m_A/\mathfrak m_A^2))$ by (\ref{equation-sequence-functorial}).
Choose a splitting
$$
\Omega_{B/\Lambda} \otimes_B k =
\Omega_{A/\Lambda} \otimes_A k \oplus K
$$
of $k$-vector space. The map $\text{d} : B \to \Omega_{B/\Lambda}$
induces via the projection onto $K$ a map $D : B \to K$. Set
$C = \{b \in B \mid D(b) = 0\}$. The Leibniz rule shows that this is
a $\Lambda$-subalgebra of $B$. Let $\overline{x} \in k$. Choose $x \in B$
mapping to $\overline{x}$. If $D(x) \not = 0$, then we can find an element
$y \in \mathfrak m_B$ such that $D(y) = D(x)$. Hence $x - y \in C$ is
an element which maps to $\overline{x}$. Thus $C \to k$ is surjective
and $C$ is an object of $\mathcal{C}_\Lambda$. Similarly, pick
$\omega \in \Im(\text{d}_A)$. We can find $x \in \mathfrak m_B$
such that $\text{d}_B(x)$ maps to $\omega$ by (1). If $D(x) \not = 0$, then
we can find an element $y \in \mathfrak m_B$ which maps to zero
in $\mathfrak m_A/\mathfrak m_A^2$ such that $D(y) = D(x)$.
Hence $z = x - y$ is an element of $\mathfrak m_C$ whose
image $\text{d}_C(z) \in \Omega_{C/k} \otimes_C k$ maps to $\omega$.
Hence $\Im(\text{d}_C) \to \Im(\text{d}_A)$ is surjective.
We conclude that $C \to A$ is surjective by (1). Hence $C \to B$ is
surjective by assumption. Hence $D = 0$, i.e., $K = 0$, i.e., (2)(c) holds.
This finishes the proof of (2).

\medskip\noindent
Proof of (3). If $k'/k$ is separable, then
$H_1(L_{k/\Lambda}) =
\mathfrak m_\Lambda/\mathfrak m_\Lambda^2 \otimes_{k'} k$, see
Lemma \ref{lemma-H1-separable-case}.
Hence $\Im(\text{d}_A) =
\mathfrak m_A/(\mathfrak m_\Lambda A + \mathfrak m_A^2)$
and similarly for $B$. Thus (3) follows from (2).

\medskip\noindent
Proof of (4). A section $s$ of $f$ is not surjective (by definition a
small extension has nontrivial kernel), hence $f$ is not essentially
surjective. Conversely, assume $f$ is a small extension but not an
essential surjection. Choose a ring map $C \to B$ in $\mathcal{C}_\Lambda$
which is not surjective, such that $C \to A$ is surjective. Let
$C' \subset B$ be the image of $C \to B$. Then $C' \not = B$ but
$C'$ surjects onto $A$. Since $f : B \to A$ is a small extension,
$\text{length}_C(B) = \text{length}_C(A) + 1$. Thus
$\text{length}_C(C') \leq \text{length}_C(A)$ since
$C'$ is a proper subring of $B$. But $C' \to A$ is surjective, so in
fact we must have $\text{length}_C(C') = \text{length}_C(A)$ and
$C' \to A$ is an isomorphism which gives us our section.
\end{proof}